\section{Reasoning 预算分配策略}

\subsection{Reasoning Compute 的权衡}

Reasoning Model 可以自主运行数小时，因此一个关键决策是：\textbf{在每个子任务上分配多少推理计算量}？

GPT-5.2-Codex 提供了四档 Reasoning 模式：Low、Medium、High、XHigh。更高的 Reasoning 预算意味着：

\begin{table}[H]
\centering
\caption{Reasoning 预算的权衡}
\begin{tabular}{p{3cm}p{5cm}p{5cm}}
\toprule
\textbf{维度} & \textbf{高 Reasoning 预算} & \textbf{低 Reasoning 预算} \\
\midrule
推理质量 & 更深入的分析，更少的逻辑错误 & 可能遗漏复杂依赖关系 \\
Token 消耗 & 可达 2 倍以上 & 更经济 \\
时间消耗 & 更慢，可能触发超时 & 更快 \\
适用阶段 & 规划、验证 & 常规实现 \\
\bottomrule
\end{tabular}
\end{table}

\subsection{Reasoning Sandwich 策略}

LangChain 实验发现了一个有趣的启发式——\textbf{Reasoning Sandwich}（推理三明治）：

\begin{importantbox}{Reasoning Sandwich: XHigh-High-XHigh}
\begin{itemize}[leftmargin=1.5em]
    \item \textbf{规划阶段}（XHigh）：充分理解问题，制定高质量计划。某些 Terminal Bench 任务非常复杂，好的计划能帮助更快到达可行解
    \item \textbf{实现阶段}（High）：按计划执行，不需要最高级别推理
    \item \textbf{验证阶段}（XHigh）：仔细检查结果，捕获错误，确保提交质量
\end{itemize}
这种"两头重、中间轻"的分配策略在实践中表现最优。
\end{importantbox}

\subsection{实验数据}

LangChain 的实验结果揭示了一个反直觉的发现：

\begin{itemize}[leftmargin=1.5em]
    \item 全程 XHigh：仅 53.9\%（大量任务因超时失败）
    \item 全程 High：63.6\%
    \item XHigh-High-XHigh Sandwich：66.5\%
\end{itemize}

全程使用最高推理预算反而表现最差，因为 Reasoning 消耗的时间导致大量任务超时。这说明\textbf{Reasoning 预算并非越高越好}，而是需要"把好钢用在刀刃上"。

\begin{knowledgebox}{Adaptive Reasoning 的未来方向}
Claude 和 Gemini 系列模型已经开始支持 Adaptive Reasoning——让模型自行决定在每个步骤投入多少推理计算量。另一种思路是多模型协作架构：用大模型做规划，将实现工作交给小模型。这些都是 Harness Engineering 的前沿探索方向。
\end{knowledgebox}

\subsection{本章小结}

Reasoning 预算分配是 Agent 性能的重要调节杠杆。全程最高预算会因超时反而降低性能。XHigh-High-XHigh 的"三明治"策略通过在规划和验证阶段投入更多推理资源，在实现阶段适度降低，实现了最优的性能-效率平衡。


